\chapter{Beschrijvende statistiek}\label{ch:g10:stats}

De statistiek vat een grote verzameling getallen --- punten, lonen,
temperaturen --- samen in enkele goed gekozen kengetallen. Dit hoofdstuk voert
de twee families kengetallen in: die welke het \emph{\omterm{prop:g10:coordgeom:midpoint}{midden}} van de gegevens
aanwijzen (\omterm{def:g10:stats:mean}{gemiddelde}, \omterm{def:g10:stats:median}{mediaan}) en die welke hun \emph{spreiding} meten
(\omterm{def:g10:stats:quartiles}{variatiebreedte}, \omterm{def:g10:stats:quartiles}{kwartielen}). Een diepere behandeling, met de
standaardafwijking, volgt in \cref{ch:g11:stat}.

\section{Gegevens en frequenties}

\begin{definition}[Statistische reeks]\label{def:g10:stats:series}
Een \emph{statistische reeks}\index{statistische reeks} is een lijst waarden
$x_1, x_2, \dots, x_N$ die op een populatie van omvang $N$ waargenomen zijn.
Komt een waarde $x$ $n$ keer voor, dan is $n$ haar \emph{aantal} en is het
quotiënt
\[
f = \frac{n}{N}
\]
haar \emph{frequentie}\index{frequentie} (vaak als percentage uitgedrukt). De
frequenties van alle waarden tellen op tot $1$.
\end{definition}

\begin{example}\label{ex:g10:stats:frequencies}
De punten van $20$ leerlingen op een toets van $5$ punten:
\begin{center}
\renewcommand{\arraystretch}{1.2}
\begin{tabular}{c|cccccc}
punten $x$ & $0$ & $1$ & $2$ & $3$ & $4$ & $5$ \\
\hline
aantal $n$ & $1$ & $2$ & $4$ & $6$ & $5$ & $2$ \\
\omterm{def:g10:stats:series}{frequentie} $f$ & $0.05$ & $0.10$ & $0.20$ & $0.30$ & $0.25$ & $0.10$ \\
\end{tabular}
\end{center}
Controle: $1 + 2 + 4 + 6 + 5 + 2 = 20$, en de \omterm{def:g10:stats:series}{frequenties} tellen op tot $1$.
\end{example}

\begin{omfigure}
\begin{tikzpicture}
\begin{axis}[
  width=8.4cm, height=5.2cm,
  ybar, bar width=14pt,
  xmin=-0.7, xmax=5.7, ymin=0, ymax=7.4,
  xtick={0,1,2,3,4,5},
  ytick={1,2,3,4,5,6},
  xlabel={punten}, ylabel={aantal},
  tick label style={font=\small},
  label style={font=\small}]
  \addplot[fill=omDef!30, draw=omDef] coordinates
    {(0,1) (1,2) (2,4) (3,6) (4,5) (5,2)};
\end{axis}
\end{tikzpicture}

{\small Het staafdiagram van de toetspunten: één staaf per waarde, met een
hoogte gelijk aan het aantal.}
\end{omfigure}

\section{Kengetallen van het midden}

\begin{definition}[Gemiddelde]\label{def:g10:stats:mean}
Het \emph{gemiddelde}\index{gemiddelde} van de reeks is
\[
\bar x = \frac{x_1 + x_2 + \dots + x_N}{N}.
\]
Komen de waarden $v_1, \dots, v_k$ voor met aantallen $n_1, \dots, n_k$, dan
bereken je hetzelfde getal als het \emph{gewogen gemiddelde}
\[
\bar x = \frac{n_1 v_1 + n_2 v_2 + \dots + n_k v_k}{n_1 + n_2 + \dots + n_k}.
\]
\end{definition}

\begin{example}\label{ex:g10:stats:mean}
Voor de toets van \cref{ex:g10:stats:frequencies}:
\[
\bar x = \frac{1 \times 0 + 2 \times 1 + 4 \times 2 + 6 \times 3
+ 5 \times 4 + 2 \times 5}{20}
= \frac{0 + 2 + 8 + 18 + 20 + 10}{20}
= \frac{58}{20} = 2.9 .
\]
\end{example}

\begin{definition}[Mediaan]\label{def:g10:stats:median}
De \emph{mediaan}\index{mediaan} is een waarde die de \emph{gesorteerde}
reeks in twee helften splitst: minstens de helft van de waarden is $\leq$ de
mediaan, en minstens de helft is $\geq$ de mediaan. In de praktijk sorteer je
de $N$ waarden;
\begin{itemize}
  \item is $N$ oneven, dan is de mediaan de middelste waarde, op plaats
        $\frac{N+1}{2}$;
  \item is $N$ even, neem dan het \omterm{prop:g10:coordgeom:midpoint}{midden} van de twee middelste waarden, op
        de plaatsen $\frac N2$ en $\frac N2 + 1$.
\end{itemize}
\end{definition}

\begin{example}\label{ex:g10:stats:median}
Een gesorteerde reeks van $7$ huizenprijzen (in duizenden): $120$, $150$,
$160$, $180$, $210$, $240$, $900$. De \omterm{def:g10:stats:median}{mediaan} is de $4$de waarde, $180$. Het
\omterm{def:g10:stats:mean}{gemiddelde} is $\frac{1960}{7} = 280$ --- groter dan $6$ van de $7$ prijzen!
Eén extreme waarde ($900$) trekt het \omterm{def:g10:stats:mean}{gemiddelde} veel harder mee dan de
\omterm{def:g10:stats:median}{mediaan}: de \omterm{def:g10:stats:median}{mediaan} is \emph{robuust}, het \omterm{def:g10:stats:mean}{gemiddelde} niet.
\end{example}

\section{Kengetallen van de spreiding}

\begin{definition}[Variatiebreedte, kwartielen]\label{def:g10:stats:quartiles}
Voor een gesorteerde reeks:
\begin{itemize}
  \item de \emph{variatiebreedte}\index{variatiebreedte} is het verschil
        tussen de grootste en de kleinste waarde;
  \item het \emph{eerste kwartiel}\index{kwartiel} $Q_1$ is de kleinste
        waarde waarvoor minstens een kwart ($25\%$) van de waarden $\leq Q_1$
        is; het \emph{derde kwartiel} $Q_3$ is de kleinste waarde waarvoor
        minstens drie kwart ($75\%$) van de waarden $\leq Q_3$ is;
  \item de \emph{interkwartielafstand}\index{interkwartielafstand}
        $Q_3 - Q_1$ meet de spreiding van de middelste helft van de gegevens.
\end{itemize}
\end{definition}

\begin{method}[De kwartielen bepalen]\label{met:g10:stats:quartiles}
Voor een reeks van $N$ gesorteerde waarden:
\begin{enumerate}
  \item bereken $\frac N4$; is dat geen \omterm{def:g10:numbers:sets}{geheel getal}, rond dan \emph{naar
        boven} af tot het volgende gehele getal; $Q_1$ is de waarde op die
        plaats;
  \item bereken $\frac{3N}{4}$; rond zo nodig naar boven af; $Q_3$ is de
        waarde op die plaats.
\end{enumerate}
\end{method}

\begin{example}\label{ex:g10:stats:quartiles}
Neem de $11$ gesorteerde waarden
\[
2,\ 3,\ 5,\ 5,\ 6,\ 7,\ 8,\ 9,\ 9,\ 10,\ 12 .
\]
\omterm{def:g10:stats:median}{Mediaan}: plaats $\frac{11+1}{2} = 6$, dus de \omterm{def:g10:stats:median}{mediaan} is $7$.
$Q_1$: $\frac{11}{4} = 2.75$, naar boven $3$; de $3$de waarde is $5$.
$Q_3$: $\frac{33}{4} = 8.25$, naar boven $9$; de $9$de waarde is $9$.
\omterm{def:g10:stats:quartiles}{Variatiebreedte}: $12 - 2 = 10$; \omterm{def:g10:stats:quartiles}{interkwartielafstand}: $9 - 5 = 4$.
\end{example}

\begin{omfigure}
\begin{tikzpicture}[scale=0.85]
  \draw[-Stealth] (1,0) -- (13.4,0) node[below, font=\small] {waarden};
  \foreach \x in {2,4,6,8,10,12}
    { \draw (\x,-0.07) -- (\x,0.07);
      \node[below, font=\footnotesize] at (\x,-0.1) {\x}; }
  \draw[omDef, thick, fill=omDef!12] (5,0.5) rectangle (9,1.3);
  \draw[omDef, very thick] (7,0.5) -- (7,1.3);
  \draw[omDef, thick] (2,0.9) -- (5,0.9);
  \draw[omDef, thick] (9,0.9) -- (12,0.9);
  \draw[omDef, thick] (2,0.7) -- (2,1.1);
  \draw[omDef, thick] (12,0.7) -- (12,1.1);
  \node[above, font=\small] at (5,1.32) {$Q_1$};
  \node[above, font=\small] at (7,1.32) {mediaan};
  \node[above, font=\small] at (9,1.32) {$Q_3$};
  \node[above, font=\small] at (2,1.12) {min};
  \node[above, font=\small] at (12,1.12) {max};
\end{tikzpicture}

{\small De boxplot van \cref{ex:g10:stats:quartiles}: de doos beslaat het
interkwartielinterval $\intcc{Q_1}{Q_3}$ en bevat de middelste helft van de
gegevens; de snorharen reiken tot de extreme waarden.}
\end{omfigure}

\begin{example}[Twee groepen vergelijken]\label{ex:g10:stats:comparison}
Twee klassen maken dezelfde toets (punten op $20$):
\begin{center}
\renewcommand{\arraystretch}{1.2}
\begin{tabular}{c|ccc}
 & \omterm{def:g10:stats:median}{mediaan} & $Q_1$ & $Q_3$ \\
\hline
klas A & $12$ & $10$ & $14$ \\
klas B & $12$ & $6$ & $17$ \\
\end{tabular}
\end{center}
Dezelfde \omterm{def:g10:stats:median}{mediaan}, heel verschillende spreiding: klas A is homogeen (de helft
van haar punten ligt in $\intcc{10}{14}$), klas B is heterogeen (haar
middelste helft loopt over $\intcc{6}{17}$). Kengetallen van het \omterm{prop:g10:coordgeom:midpoint}{midden}
vertellen alleen nooit het hele verhaal.
\end{example}

\begin{remark}
Noch het \omterm{def:g10:stats:mean}{gemiddelde} noch de \omterm{def:g10:stats:median}{mediaan} is ``beter'': het \omterm{def:g10:stats:mean}{gemiddelde} gebruikt
elke waarde (en is nodig om totalen te berekenen), de \omterm{def:g10:stats:median}{mediaan} weerstaat
extreme waarden. Een ernstige samenvatting van een reeks geeft minstens één
kengetal van het \omterm{prop:g10:coordgeom:midpoint}{midden} \emph{en} één van de spreiding.
\end{remark}

\section{Oefeningen}

\begin{exercise}[$\star$]\label{exo:g10:stats:1}
Hier staan de aantallen boeken die $15$ leerlingen in een jaar lazen:
\[
0,\ 1,\ 1,\ 2,\ 2,\ 2,\ 3,\ 3,\ 4,\ 4,\ 5,\ 6,\ 7,\ 8,\ 12 .
\]
Bereken het \omterm{def:g10:stats:mean}{gemiddelde}, de \omterm{def:g10:stats:median}{mediaan} en de \omterm{def:g10:stats:quartiles}{variatiebreedte} van deze reeks.
\end{exercise}

\begin{exercise}[$\star$]\label{exo:g10:stats:2}
Er werd $50$ keer met een dobbelsteen gegooid:
\begin{center}
\renewcommand{\arraystretch}{1.2}
\begin{tabular}{c|cccccc}
uitkomst & $1$ & $2$ & $3$ & $4$ & $5$ & $6$ \\
\hline
aantal & $7$ & $9$ & $8$ & $10$ & $6$ & $10$ \\
\end{tabular}
\end{center}
Bereken de \omterm{def:g10:stats:series}{frequentie} van elke uitkomst en de \omterm{def:g10:stats:mean}{gemiddelde} uitkomst.
\end{exercise}

\begin{exercise}[$\star$]\label{exo:g10:stats:3}
Bepaal de \omterm{def:g10:stats:median}{mediaan}, $Q_1$ en $Q_3$ van de gesorteerde reeks
\[
1,\ 2,\ 4,\ 4,\ 5,\ 6,\ 6,\ 7,\ 8,\ 9,\ 9,\ 11 \qquad (N = 12),
\]
en teken de bijhorende boxplot.
\end{exercise}

\begin{exercise}[$\star$]\label{exo:g10:stats:4}
Het \omterm{def:g10:stats:mean}{gemiddelde} van $24$ toetsresultaten is $11.5$. Een $25$ste leerling maakt
de toets en haalt $19$. Wat is het nieuwe klasgemiddelde?
\end{exercise}

\begin{exercise}[$\star\star$]\label{exo:g10:stats:5}
Een leerling heeft na $5$ toetsen, alle even zwaar, een \omterm{def:g10:stats:mean}{gemiddelde} van $12$.
Welk resultaat op de $6$de toets zou het \omterm{def:g10:stats:mean}{gemiddelde} naar $13$ tillen? Is een
\omterm{def:g10:stats:mean}{gemiddelde} van $14$ haalbaar (de punten staan op $20$)?
\end{exercise}

\begin{exercise}[$\star\star$]\label{exo:g10:stats:6}
Verzin een reeks van $7$ niet-negatieve waarden met \omterm{def:g10:stats:median}{mediaan} $10$ en een
\omterm{def:g10:stats:mean}{gemiddelde} groter dan $20$; en daarna een met \omterm{def:g10:stats:median}{mediaan} $10$ en een \omterm{def:g10:stats:mean}{gemiddelde}
kleiner dan $6$. Waarom kan het \omterm{def:g10:stats:mean}{gemiddelde} van zo'n reeks nooit kleiner zijn
dan $\frac{40}{7}$? Wat tonen deze voorbeelden?
\end{exercise}

\begin{exercise}[$\star\star$]\label{exo:g10:stats:7}
In een bedrijf verdienen de $9$ werknemers elk $1800$ per maand en de
directeur $12000$.
\begin{enumerate}
  \item Bereken het \omterm{def:g10:stats:mean}{gemiddelde} en het mediane loon.
  \item Welk kengetal beschrijft hier het best een ``typisch'' loon? Waarom?
\end{enumerate}
\end{exercise}

\begin{exercise}[$\star\star$]\label{exo:g10:stats:8}
Een klas met $12$ jongens heeft een \omterm{def:g10:stats:mean}{gemiddelde} lengte van $168$ cm; de $18$
meisjes van dezelfde klas hebben een \omterm{def:g10:stats:mean}{gemiddelde} lengte van $163$ cm. Bereken
de \omterm{def:g10:stats:mean}{gemiddelde} lengte van de hele klas. (Let op: dat is \emph{niet} het \omterm{prop:g10:coordgeom:midpoint}{midden}
van $168$ en $163$ --- weeg de \omterm{def:g10:stats:mean}{gemiddelden} met de groepsgroottes.)
\end{exercise}

\begin{exercise}[$\star\star\star$]\label{exo:g10:stats:9}
Een reeks van $N$ waarden $x_1, \dots, x_N$ heeft \omterm{def:g10:stats:mean}{gemiddelde} $\bar x$. Elke
waarde wordt omgezet in $y_i = a x_i + b$ voor vaste getallen $a$ en $b$.
\begin{enumerate}
  \item Toon aan dat het \omterm{def:g10:stats:mean}{gemiddelde} van de nieuwe reeks $a \bar x + b$ is.
  \item De temperaturen van een week hebben, in graden Celsius, \omterm{def:g10:stats:mean}{gemiddelde}
        $20$. Wat is hun \omterm{def:g10:stats:mean}{gemiddelde} in graden Fahrenheit ($F = 1.8\,C + 32$)?
\end{enumerate}
\end{exercise}

\section{Opgave: de gemiddelde mens bestaat niet}

\begin{problem}[{Weekendopgave --- gemiddelde tegen mediaan: uitschieters,
puntenaanpassingen, het magazijn langs de weg, en de cockpit die voor
niemand ontworpen was}]
\label{pb:g10:stats:1}
In 1950 mat de Amerikaanse luchtmacht $4\,000$ piloten op tien
lichaamsmaten en bouwde ze de cockpit voor de \emph{\omterm{def:g10:stats:mean}{gemiddelde} man}. Hij
paste --- zoals een luitenant al tellend ontdekte --- vrijwel niemand. De
statistiek vat menigten samen in enkele getallen, en elke samenvatting
vertelt de waarheid over het ene en liegt over het andere. Deze opgave
daagt het \omterm{def:g10:stats:mean}{gemiddelde} en de \omterm{def:g10:stats:median}{mediaan} (\cref{def:g10:stats:mean},
\cref{def:g10:stats:median}) voor de rechtbank, laat elk tonen waar het het
best in is, en sluit de zaak van de cockpit.

\medskip
\noindent\textbf{Deel I --- Twee samenvattingen voor de rechter.}
\begin{enumerate}
  \item In een straat worden zeven huizen verkocht, in duizenden euro's:
        $200$, $210$, $220$, $230$, $240$, $250$ --- en één landhuis van
        $1\,200$. Bereken de \omterm{def:g10:stats:mean}{gemiddelde} en de mediane prijs. Welk getal
        beschrijft ``een huis in deze straat''?
  \item Haal het landhuis weg en bereken beide opnieuw. Hoeveel schoof elke
        samenvatting op? Formuleer de moraal over uitschieters.
  \item Een volkstellingstabel: van $100$ gezinnen hebben er $20$ geen kind,
        $25$ één, $30$ twee, $15$ drie en $10$ vier. Bereken het \omterm{def:g10:stats:mean}{gemiddelde}
        aantal kinderen. Geen enkel gezin ``heeft $1.7$ kinderen'' --- welke
        werkelijke grootheid legt het \omterm{def:g10:stats:mean}{gemiddelde} dan vast? (Vermenigvuldig het
        met $100$.)
  \item Zoek voor de vijftien punten
        $4, 6, 7, 8, 9, 10, 10, 11, 12, 12, 13, 14, 15, 17, 19$ de \omterm{def:g10:stats:median}{mediaan} en
        de \omterm{def:g10:stats:quartiles}{kwartielen} $Q_1$, $Q_3$ (\cref{met:g10:stats:quartiles}).
  \item Geef voor dezelfde reeks de \omterm{def:g10:stats:quartiles}{variatiebreedte} en de
        \omterm{def:g10:stats:quartiles}{interkwartielafstand}. Wat meet elk van beide?
\end{enumerate}

\medskip
\noindent\textbf{Deel II --- Punten aanpassen en klassen samenvoegen.}
\begin{enumerate}[resume]
  \item Uit \cref{exo:g10:stats:9}: elke waarde met $+b$ verschuiven
        verschuift \omterm{def:g10:stats:mean}{gemiddelde}, \omterm{def:g10:stats:median}{mediaan} en \omterm{def:g10:stats:quartiles}{kwartielen} met $+b$; met $a$
        schalen schaalt ze met $a$. Wat gebeurt er met de \omterm{def:g10:stats:quartiles}{variatiebreedte} en
        de \omterm{def:g10:stats:quartiles}{interkwartielafstand} bij een verschuiving? En bij een schaling?
        Verantwoord.
  \item Een toets heeft \omterm{def:g10:stats:mean}{gemiddelde} $9$ (op $20$) en de leerkracht wil
        \omterm{def:g10:stats:mean}{gemiddelde} $12$. Er worden twee aanpassingen voorgesteld: iedereen
        $3$ punten erbij, of elk resultaat met $\frac43$ vermenigvuldigen.
        Bereken het lot van een leerling met $6$ en van een met $18$ onder
        elke aanpassing. Welke aanpassing begunstigt wie --- en welke stoot
        haar hoofd tegen het plafond?
  \item Klas A ($20$ leerlingen) heeft \omterm{def:g10:stats:mean}{gemiddelde} $12$; klas B ($30$
        leerlingen) heeft \omterm{def:g10:stats:mean}{gemiddelde} $10$. Bereken het \omterm{def:g10:stats:mean}{gemiddelde} van de $50$
        leerlingen samen, en leg uit waarom het geen $11$ is (de les over
        gewogen \omterm{def:g10:stats:mean}{gemiddelden} uit de weekendopgave over het harmonisch
        \omterm{def:g10:stats:mean}{gemiddelde} in het onderbouwvolume, nu in de klas).
  \item Toon aan dat medianen dit spel weigeren: bereken de medianen van
        $A = \{1, 2, 3\}$ en $B = \{0, 10\}$, daarna de \omterm{def:g10:stats:median}{mediaan} van de
        samengevoegde lijst, en vergelijk met het naar omvang gewogen
        \omterm{def:g10:stats:mean}{gemiddelde} van de twee medianen.
  \item Kunstschaatsen: vijf juryleden geven $5.2$, $5.3$, $5.4$, $5.5$,
        $9.9$. Bereken het \omterm{def:g10:stats:mean}{gemiddelde}, de \omterm{def:g10:stats:median}{mediaan} en het \emph{getrimde
        \omterm{def:g10:stats:mean}{gemiddelde}} (laat de laagste en de hoogste vallen en middel de rest).
        Waarom trimmen de meeste beoordeelde sporten?
\end{enumerate}

\medskip
\noindent\textbf{Deel III --- Verdelingen lezen.}
\begin{enumerate}[resume]
  \item Twee fabrieken maken allebei bouten met een \omterm{def:g10:stats:mean}{gemiddelde} lengte van
        $50$ mm. Fabriek A: $Q_1 = 49.9$, $Q_3 = 50.1$; fabriek B:
        $Q_1 = 48$, $Q_3 = 52$. Je hebt onderling verwisselbare bouten nodig:
        welke fabriek, en welk kengetal besliste?
  \item In een kliniek is de mediane wachttijd $10$ minuten, maar de
        \omterm{def:g10:stats:mean}{gemiddelde} wachttijd $25$. Welke vorm van de gegevens verraadt dat
        paar? Maak een reeks van vijf waarden met precies deze \omterm{def:g10:stats:median}{mediaan} en dit
        \omterm{def:g10:stats:mean}{gemiddelde}.
  \item ``Je baby zit op het 75ste percentiel voor lengte'' --- wat betekent
        die zin? En kan elk kind van een stad ``boven het \omterm{def:g10:stats:mean}{gemiddelde}'' zijn?
        Boven de \omterm{def:g10:stats:median}{mediaan}? Leg het verschil uit.
  \item Klas A: min $4$, $Q_1\,8$, \omterm{def:g10:stats:median}{mediaan} $11$, $Q_3\,13$, max $17$. Klas
        B: min $2$, $Q_1\,9$, \omterm{def:g10:stats:median}{mediaan} $11$, $Q_3\,12$, max $19$. Vergelijk:
        typisch punt, homogeniteit van de middelste helft, wildheid van de
        uitersten.
  \item Het lijstje van de journalist: een kop zegt ``gemiddeld loon bij
        TechCorp: $5\,000$ euro''. Noem de drie vragen die deze opgave je
        geleerd heeft te stellen voordat je gelooft dat de kop iets zegt over
        een typische werknemer.
\end{enumerate}

\medskip
\noindent\textbf{Deel IV --- De zaak van de cockpit.}
\begin{enumerate}[resume]
  \item Maak je eigen bewijsstuk: verzin vijf huizenprijzen met \omterm{def:g10:stats:mean}{gemiddelde}
        $300$ en \omterm{def:g10:stats:median}{mediaan} $200$ (in duizenden euro's).
  \item Het \omterm{def:g10:stats:mean}{gemiddelde} is het evenwichtspunt: bewijs uit de definitie dat de
        afwijkingen van het \omterm{def:g10:stats:mean}{gemiddelde} elkaar altijd opheffen,
        $\sum (x_i - \bar x) = 0$, en ga het na op je reeks uit vraag 16.
  \item De \omterm{def:g10:stats:median}{mediaan} bespaart boodschappen: langs een weg staan winkels op
        kilometer $1$, $3$, $4$, $9$ en $13$, en een magazijn moet de
        \emph{totale afstand} tot de vijf winkels zo klein mogelijk maken.
        Bereken die totale afstand voor een magazijn op de \omterm{def:g10:stats:median}{mediaan} ($4$), op
        het \omterm{def:g10:stats:mean}{gemiddelde} ($6$) en op $5$. Wie wint? (Dit geldt algemeen: de
        \omterm{def:g10:stats:median}{mediaan} maakt de totale absolute afstand minimaal --- toets gerust een
        andere plek.)
  \item De cockpit, opgelost: stel dat een piloot $30\,\%$ kans heeft om op
        één lichaamsmaat ``dicht bij het \omterm{def:g10:stats:mean}{gemiddelde}'' te zitten, onafhankelijk
        over $10$ maten. Schat de kans om op \emph{alle tien} dicht bij het
        \omterm{def:g10:stats:mean}{gemiddelde} te zitten --- en breng het resultaat in overeenstemming met
        wat de luitenant bij $4\,000$ piloten vond. Wat bouwde de luchtmacht
        in plaats van de \omterm{def:g10:stats:mean}{gemiddelde} cockpit?
  \item Slotstuk --- de handleiding, telkens één regel: gaat het om totalen
        (een begroting, een oogst), gebruik dan het \dots; gaat het om het
        typische individu (lonen, wachttijden), gebruik dan de \dots; gaat het
        om eerlijkheid van de spreiding, meld dan de \dots; dreigt er
        besmetting, gebruik dan het \dots. Sluit de zaak met de titel van deze
        opgave.
\end{enumerate}
\end{problem}
